\documentclass[a4paper,12pt]{article}

\usepackage[utf8]{inputenc}
\usepackage[T2A]{fontenc}
\usepackage[russian]{babel}
\usepackage{geometry}
\geometry{top=2cm, bottom=2cm, left=3cm, right=1.5cm}
\usepackage{hyperref}
\usepackage{graphicx}
\usepackage{enumitem}
\usepackage{indentfirst}

\begin{document}


\begin{titlepage}
    \begin{center}
        \vspace*{5cm}
        
        \textbf{\Large ОТЧЕТ ПО ПРОЕКТУ}
        
        \vspace{2cm}
        
        {\Large \textbf{Интерактивная карта: Ростовская область в годы Великой Отечественной войны (1941--1943)}}
        
        \vspace{4cm}
        
        \begin{flushright}
            \textbf{Участники проекта:}\\
            Лазимов Д.И. (Тимлид) \\
            Артеменко А.О. \\
            Бармакова К.C. \\ 
        \end{flushright}
        
        \vfill
        
        {\large 2026} 
    \end{center}
\end{titlepage}

\newpage

\section*{Синопсис проекта}
«Ростовская область в годы Великой Отечественной войны (1941--1943)» — это масштабный цифровой интерактивный проект, представляющий собой веб-приложение для изучения локальной истории. В центре проекта — интерактивная географическая карта, на которой визуализирована хронология боевых действий, периодов оккупации и поэтапного освобождения населенных пунктов Дона. Проект охватывает период с 22 июня 1941 года по конец августа 1943 года. 

Ключевая особенность приложения — динамический таймлайн со встроенным плеером, который позволяет пользователю наблюдать за развитием событий в режиме ускоренного реального времени. Приложение обеспечивает глубокое погружение в материал: плеер автоматически останавливает воспроизведение, когда на карте появляется новое историческое событие, и выводит подробную справку в боковую панель. Все отметки категоризированы (битвы, оккупация, освобождение, трагедии Холокоста), а для обеспечения контекста на экране отображается название текущей глобальной военной операции.

\section*{Целевая аудитория проекта и актуальность для нее}
\begin{enumerate}
    \item \textbf{Школьники и студенты} – проект послужит наглядным и интерактивным образовательным пособием. Визуализация сложного исторического материала с геопривязкой помогает лучше усвоить хронологию и географию событий, заменяя скучные таблицы динамичным интерактивом.
    \item \textbf{Учителя истории, краеведы и преподаватели} – получат современный цифровой инструмент для проведения уроков и лекций. Карту можно выводить на проектор для пошагового разбора таких событий, как прорыв Миус-фронта, оборона Ростова или Тацинский рейд.
    \item \textbf{Жители Ростовской области, энтузиасты и туристы} – смогут открыть для себя новые исторические факты о родном крае. Приложение помогает осознать масштаб происходившего буквально на соседних улицах. Например, героическая оборона вокзала батальоном Мадояна или трагедия в Змиевской балке.
    \item \textbf{Экскурсоводы} – проект может служить надежным цифровым справочником для создания новых краеведческих маршрутов по местам боевой славы региона.
\end{enumerate}

\section*{Описание процесса работы над исторической частью проекта}
Сбор, систематизация и оцифровка исторического материала проходили в несколько этапов:
\begin{itemize}
    \item \textbf{Сбор данных:} Изучение открытых источников, исторических справок, энциклопедий и краеведческих материалов, описывающих ход ВОВ на территории Ростовской области.
    \item \textbf{Фильтрация и расширение базы:} Был сформирован список из более чем 30 ключевых событий. В него вошли не только крупнейшие вехи (освобождение Ростова, падение Таганрога), но и важнейшие локальные подвиги. Например, оборона Зеленого острова, боевые действия в Вёшенской, героическая оборона ж/д вокзала.
    \item \textbf{Геопривязка:} Каждому событию были присвоены точные GPS-координаты с использованием открытых картографических сервисов для корректного отображения на слое Leaflet. Для масштабных локальных событий. Например, зона боев батальона Мадояна высчитывались полигоны для заливки областей.
    \item \textbf{Структурирование и копирайтинг:} События были классифицированы по типам: battle, occupation, liberation, holocaust. Для каждого была написана информативная справка, адаптированная для удобного чтения с цифровых устройств. Всем элементам присвоены точные даты для корректной интеграции с логикой плеера.
    \item \textbf{Макроконтекст:} Выделены основные периоды масштабных операций (например, "Донбасская наступательная операция", "Ростовская операция в 1943"), чтобы пользователь всегда понимал общий контекст происходящего на карте.
\end{itemize}

\section*{Описание процесса работы над технической частью проекта}
Проект реализован в архитектуре Single Page Application с использованием современного фронтенд-стека.

\textbf{Технические компетенции и стек:}
\begin{itemize}
    \item \textbf{React 18 (TypeScript)} и сборщик \textbf{Vite} — для структурированной и быстрой разработки компонентов.
    \item \textbf{Tailwind CSS} — для создания современного и адаптивного пользовательского интерфейса (Dark Mode).
    \item \textbf{React-Leaflet (Leaflet.js)} — картографический движок. Применены кастомные маркеры (DivIcon) для цветового кодирования типов событий, добавлены слои полигонов.
    \item \textbf{Lucide-React} и \textbf{date-fns} — для иконок и локализованного форматирования дат соответственно.
\end{itemize}

\textbf{Алгоритм действий и ключевые технические решения:}
\begin{enumerate}
    \item Разработка структуры данных в TypeScript, выделение констант дат для ограничения шкалы времени.
    \item Интеграция интерактивной карты со стилизованной темной подложжкой (\textit{CartoDB Dark Matter}) для контрастного отображения маркеров и фронтовых линий.
    \item Реализация сложной логики таймлайна requestAnimationFrame, обеспечивающей плавный, зависимый от прошедшего реального времени, пересчет исторического времени.
    \item \textbf{Система «Умной паузы»:} Разработан алгоритм, отслеживающий появление на карте новых событий с помощью структуры данных Set. Как только историческое время совпадает с датой нового события, плеер ставится на паузу, а в боковой панели открывается детальная справка. Была успешно решена проблема сброса кэша прочитанных событий при ручной перемотке времени назад или перезапуске анимации.
    \item \textbf{Механика "затухания" событий:} Чтобы не перегружать интерфейс, была внедрена логика, при которой отметки событий автоматически пропадают с карты спустя одну виртуальную неделю после их начала. 
\end{enumerate}

\section*{Оценка результатов}
\textbf{Что получилось?} 
Был создан полноценный, интерактивный MVP исторического веб-приложения. Карта работает плавно, таймлайн корректно переводит время, а главная фича — умная автопауза — хорошо отрабатывает, позволяя успевать читать объемные исторические справки, не теряя нить повествования. Успешно визуализирована плотность событий 1941-1943 годов на территории области.

\textbf{Что можно было бы улучшить / продолжить (Точки роста):}
\begin{itemize}
    \item Добавление медиа: интеграция подлинных архивных фотографий и аудиосводок Совинформбюро для каждого события.
    \item Внедрение слоев исторических топографических карт поверх современной спутниковой/векторной карты для еще большего реализма. Например РККА 1940-х годов
    \item Адаптация под мобильные устройства: переработка интерфейса для меньших экранов.
\end{itemize}

\textbf{Ссылка на продукт:} \url{https://zelell1.github.io/rostov_deoccupation_interactive_map/}

\section*{Использованные источники и литература}
\begin{itemize}
    \item \textbf{Электронный банк документов «Память народа»} (\url{https://pamyat-naroda.ru/}) — использовался для уточнения точной хронологии военных операций (Ростовская наступательная, Донбасская) и движения линии фронта.
    \item \textbf{Проект «Донской временник»} (Донская государственная публичная библиотека, \url{http://www.donvrem.dspl.ru/}) — краеведческие материалы и статьи о периоде оккупации и поэтапном освобождении населенных пунктов региона.
    \item \textbf{Просветительские порталы центра «Холокост» и Мемориала «Змиевская балка»} — исторические справки о трагедии мирного населения в Ростове-на-Дону и Таганроге (Петрушина балка).
    \item \textbf{Историческая литература:} Афанасенко В.И., Кринко Е.Ф. «56-я армия в боях за Ростов. Первая победа Красной армии» — источник для детального описания формирования Ростовского стрелкового полка народного ополчения, подвига батальона Мадояна и бойцов НКВД.
\end{itemize}

\section*{Вклад участников}

\begin{itemize}
\item  Лазимов Д.И. - 34 \% 
\item  Артеменко А.О. - 33 \% 
\item  Бармакова К.C.- 33 \%
\end{itemize}

\end{document}
